\section{What Modern Genomic Evidence Can Identify}
\label{sec:genomic-evidence}

Modern genomic evidence can do more than supply additional correlations
consistent with a fitted kinship model: it can help estimate components
that kinship resemblance alone does not separately identify, and test
restrictions imposed on them. Gene--environment covariance and dominance
variance illustrate why this evidence warrants different judgments about
different simplifying assumptions. Here gene--environment covariance
means correlation between genetic and environmental contributions, not
an interaction in which an effect changes with the environment.

\subsection{Random Segregation Supplies New Identifying Variation}

The decisive advance is that genotypes reveal random transmission within
families, rather than merely average resemblance between types of relatives.
Conditional on the parents' genotypes, a child's inherited alleles vary
through Mendelian segregation. Children can consequently differ in their
inherited endowments without those differences being assigned by parental
wealth, schooling, or social position. This is precisely the variation
missing when one tries to separate genetic and environmental transmission
using only average parent--child and sibling correlations.

A simple moment condition makes the advantage explicit. Let $X$ be a column
vector of offspring allele counts, $W$ the parental genotypes, and
$T=X-\mathbb E[X\mid W]$ the transmission deviation. Consider the correctly
specified additive model
\[
 Y=X^{\mathsf T}\beta+f(W)+u,\qquad \mathbb E[T u]=0.
\]
Here $\beta$ contains direct effects, while $f(W)$ allows parental-genotype
associations with family background. By construction,
$\mathbb E[T\mid W]=0$, so $\mathbb E[T f(W)]=0$ and
\[
 \mathbb E[T Y]=\mathbb E[T T^{\mathsf T}]\beta.
\]
Where the transmission covariance matrix has sufficient rank, these
moments identify the corresponding effects without requiring
$\operatorname{Cov}(X,f(W))=0$ in the population. Correlated loci can limit
which combinations are separately identified. The substantive condition
$\mathbb E[T u]=0$ still requires justification: parental conditioning
removes family-background association with transmission, but does not
license omission of relevant transmitted variants or correlated sibling
effects. Selection can also compromise the design. Nor does random
transmission require environments to be unresponsive to the child's
inherited characteristics; those responses can mediate its direct effect
\citep[Sections~3.1, 4.4--4.5]{BenjaminEtAl2026}.

Segregation also produces random variation in how much genetic material
siblings and other relatives actually share. Genomic relatedness designs
use that variation to estimate genetic variance separately from the average
resemblance associated with a family relationship
\citep{YoungEtAl2018}; see also \citet[Section~4.7]{BenjaminEtAl2026}. A PGI's segregation deviation similarly
provides a source of within-family variation. These are distinct from using
its between-family association as a causal effect. Assortment must be
accounted for when translating segregation-based estimates into population
genetic variance; the two variances are not generally identical
\citep{Young2023}.

\subsection{Evidence Already Challenges the Population-PGI Interpretation}
Population-based GWAS coefficients can combine direct genetic effects with
indirect effects, population structure, and other gene--environment
association \citep{YoungEtAl2019,YoungEtAl2022}. This is an empirical problem,
not merely a hypothetical possibility. % CITATION CHECK: Add Howe et al. (2022) primary citation after exact author-list/source verification; currently attributed through Benjamin et al.
The evidence reviewed in
\citet[Section~4.6]{BenjaminEtAl2026} includes smaller average family-based
effects for educational attainment and changes in cross-trait genetic
correlations. Treating population coefficients as unbiased direct effects
is therefore untenable as a general premise. This does not establish bias
in every coefficient or quantify the error in every claim Clark makes.

The implication for $m$ is immediate. A spouse correlation in a score
weighted by population coefficients need not measure assortment on the
direct genetic component posited by the model. Family-based weights targeting
direct effects are the appropriate starting point for this purpose and
for the book's other PGI tests insofar as they seek to distinguish direct
genetic differences from environmental explanations, including its
geographical comparisons. \citet[Section~6.4]{BenjaminEtAl2026} explicitly
recommend sufficiently well-powered family-based PGIs for assortment
studies. Better population prediction does not establish the causal
interpretation: the difference concerns the estimand as well as precision.

Family-based weights do not automatically recover $m$. Finite precision,
incomplete genetic coverage, and the particular family design's assumptions
still matter; a score correlation is not the correlation in the complete
latent endowment. These measurement issues require explicit treatment,
not interchangeability of population and family-based coefficients.

Clark discusses the non-transmitted-allele design of \citet{KongEtAl2018}
and argues that assortment and incomplete measurement can account for its
findings \citep[pp.~83--84]{Clark2026}. The criticism is not that he ignores
this design. It is that an alternative explanation of one genetic-nurture
result does not resolve the wider evidence of confounding in population
associations. His causal use of those associations requires identifying
conditions that their predictive success does not establish.

\subsection{Gene--Environment Covariance and Dominance Warrant Different Conclusions}

For gene--environment covariance, family-based designs use variation in
transmission conditional on parental genotypes to separate direct genetic
effects from associations generated by family background. Joint analyses
of offspring and parental genotypes can additionally estimate indirect
parental effects under explicit identifying assumptions. Comparisons with
population estimates therefore provide evidence about the environmental
and indirect-genetic contributions that population associations can
absorb. The evidence reviewed by \citet[Sections~4.3--4.7]{BenjaminEtAl2026}
includes substantial gene--environment correlation for educational
attainment. It challenges the exclusion of such covariance from the model;
it does not, without further assumptions, identify its entire magnitude
or allocation across kinships. In particular, subtracting family-based
from population-based coefficients does not automatically recover a single
$\operatorname{Cov}(G,U)$: the difference can reflect several sources of
confounding, and the mapping depends on the genetic coverage and the
family design. Environmental responses caused by the child's own genotype
also remain part of the direct effect, as defined in Section~\ref{sec:model-components}.

The evidence on dominance supports a different conclusion. \citet{PazokitoroudiEtAl2021} directly estimate dominance contributions across a broad set of complex traits. Dominance
variance measures departures from additivity between the two alleles at
a locus. The genomic studies reviewed by
\citet[Sections~3.2 and~4.7]{BenjaminEtAl2026}, % CITATION CHECK: Okbay et al. (2022) original/supplement is not in local library; the following EA-specific statement is explicitly attributed to the JEL review.
including the educational
attainment results of Okbay and colleagues, support a small contribution
from dominance variance. Clark's decision to omit it is consequently a
reasonable approximation for educational attainment, subject to the
variants, populations, and precision covered by those studies. This is
not a finding that every form of nonadditivity is absent, nor does it
establish that dominance is negligible for every historical outcome.
More importantly, evidence supporting negligible dominance supplies no
justification for excluding gene--environment covariance or shared
environmental influences. Modern evidence can thus support Clark's
additive approximation while challenging the environmental exclusions
on which his causal interpretation rests.

\subsection{From Fitting Kinship Correlations to Testing Their Explanation}

Better family-based PGIs can help estimate $m$ independently of the
kinship correlations the model seeks to explain. With appropriate allowance
for incomplete genetic measurement, estimates of $m$, $h^2$, and spouse
phenotypic resemblance can test the contrasting restrictions of phenotypic
and genetic assortment developed above. Cross-trait direct genetic
covariances can test a strong single-endowment model; genomic signatures
of assortment and comparisons across cohorts can probe equilibrium and
stability assumptions \citep{YengoEtAl2018}; see also \citet[Sections~4.7 and~6.4]{BenjaminEtAl2026}.
The aim should be to estimate components from transmission and genomic
information, then predict kinship correlations not used in that estimation.
This makes the data informative about which model is appropriate rather
than merely which parameter values best fit a chosen model.

The lesson of the great IQ debate was that good fit to kinship correlations
does not establish the assumptions that give fitted parameters their causal
meaning \citep{Loehlin1978,Goldberger1978}. Modern genomic data offer a way
forward by supplying independent evidence about genetic assortment, direct
and indirect genetic effects, dominance, and the consequences of mating
across generations. These developments can support some of Clark's
simplifications while challenging others. Their value lies in making
consequential assumptions testable, rather than allowing another
well-fitting model to stand as evidence of its own validity.
